\documentclass[addpoints]{exam}
% \documentclass[addpoints,12pt]{exam}

\usepackage[slovene]{babel}
\usepackage{amsfonts}
\usepackage{amssymb}
\usepackage{amsmath}
\usepackage{float}
\usepackage{graphicx}
\usepackage{enumitem}
\usepackage{color}
\usepackage{eurosym}


%Komentar naslednje vrstice glede na verzijo testa -- rešitve ali prostor za odgovore:
\printanswers

% \linespread{2}

\pointsinrightmargin
\marginbonuspointname{}


\renewcommand{\solutiontitle}{\noindent\textbf{Rešitve in točkovanje:}\par\noindent}

\colorgrids
\definecolor{GridColor}{gray}{0.7}
\setlength{\gridsize}{7mm}

\extrawidth{5mm}
\setlength{\rightpointsmargin}{1.5cm}

\begin{document}

\runningfooter{}{}{}

\begin{center}
    \fbox{\fbox{\parbox{15cm}{\centering
    Prvo pisno ocenjevanja znanja \\ 1. letnik -- test C \\ 23. oktober 2024 \\ (Osnove logike in teorije množic, Naravna in cela števila, Potence)}}}
\end{center}

\vspace{0.1in}
\makebox[\textwidth]{Ime in priimek, oddelek:\enspace\hrulefill}


\begin{center}
    \chqword{Naloga}
    \chpword{Možne točke}
    \chbpword{Dodatne točke}
    \chsword{Dosežene točke}
    \chtword{Skupaj}  
    \combinedpointtable[h][questions]
    % \combinedgradetable[h][questions]
\end{center}

\vspace{0.1in}


\begin{questions}

    % 1. vprašanje
    
    \question[4]
    Določite pravilnost izjave $(A\Rightarrow  B)\land (\lnot A\lor \lnot B)$ glede na izbor izjav $A$ in $B$.

    \begin{solution}[8cm]
        \begin{table}[H]
            \begin{tabular}{ccccccc}
             $A$ & $B$ & $\lnot A$ & $\lnot B$ & $(A\Rightarrow B)$ & $(\lnot A\lor \lnot B)$ & $(A\Rightarrow B)\land(\lnot A\lor \lnot B)$\\ \hline
             P & P & N & N & P & N & N \\ 
             P & N & N & P & N & P & N \\
             N & P & P & N & P & P & P \\
             N & N & P & P & P & P & P
            \end{tabular}
        \end{table}
        Za vsako pravilno izpolnjeno vrstico, ki določa pravilnost, po $1$ točka.
    \end{solution}




    % 2. vprašanje
    \question[3]
    Določite resničnostno vrednost izjav:
    \begin{parts}
        \part 
        Število $21$ je praštevilo ali število $15$ deli število $20$.\\
        \part 
        Če število $5$ ni praštevilo, potem je praštevil končno mnogo.\\
        \part
        Ni res, da je število $8$ večje od števila $2$ ali ostanek pri deljenju s $4$ je lahko samo $1$, $2$ ali $3$.\\
    \end{parts}

    \begin{solution}
        (a): N (prva izjava je nepravilna, druga je nepravilna) \\
        (b): P (prva izjava je nepravilna, druga je nepravilna) \\
        (c): N (negirana izjava je pravilna /prva izjava je pravilna, druga izjava je nepravilna/) \\ \\
        Za vsako pravilno določeno logično vrednost po $1$ točka.
    \end{solution}



    \newpage

    % 3. vprašanje
    \question
    Dani sta množici $\mathcal{A}=\left\{0, 2, 3, 6\right\}$ in $\mathcal{B}=\left\{0, 3, 5, 8\right\}$.
    \begin{parts}
        \part[2]
        Ali velja $\mathcal{A}\subseteq\mathcal{B}$? Poiščite največjo možno množico $\mathcal{C}$, da bosta množici $\mathcal{A}$ in $\mathcal{B}$ njeni nadmnožici.
        \part[3]
        Zapišite $\mathcal{PA}$. Izračunajte koliko elementov vsebuje.
        \part[2]
        Zapišite in grafično predstavite kartezični produkt $\mathcal{A}\times\mathcal{B}$.
    \end{parts}

    % \begin{description}
    %     \item[a] $(x-a)^2+(y-b)^2-c^2=0$ 
    %     \item[b] $d^2(x-a)^2+f^2(y-b)^2-(df)^2=0$
    %     \item[c] $S(a,b)$
    %     \item[d] $e^2=d^2-f^2;\ d>f$
    % \end{description}
    % Krožnica: \fillin[a c][2cm]
    % \\ Elipsa:  \fillin[b c d][2cm]

    \begin{solution}[9cm]
        (a): Ne drži, da $\mathcal{A}\subseteq\mathcal{B}$, ker $2\in\mathcal{A}\land 2\notin\mathcal{B}$.; $\mathcal{C}=\left\{0,3\right\}$ \\
        (b): $\mathcal{PA}=\left\{\emptyset, \{0\}, \{2\}, \{3\}, \{6\}, \{0,2\}, \{0,3\}, \{0,6\}, \{2,3\}, \{2,6\}, \{3,6\}, \{0,2,3\}, \{0,2,6\}, \{0,3,6\}, \{2,3,6\}, \mathcal{A}\right\}$; $m(\mathcal{PA})=2^{m(\mathcal{A})}=2^4=16$. \\
        (c): $\mathcal{A}\times\mathcal{B}=\left\{(0,0), (0,3), (0,5), (0,8), (2,0), (2,3), (2,5), (2,8), (3,0), (3,3), (3,5), (3,8), (6,0), (6,3), (6,5), (6,8)\right\}$ + grafična upodobitev \\ \\
        Pri (a) za vsak odgovor po $1$ točka. Pri (b) za zapis celotne $\mathcal{PA}$ $2$ točki, za zapis zgolj podmnožic $1$ točka; za izračun po formuli za moč potenčne množice $1$ točka. Pri (c) za zapis $\mathcal{A}\times\mathcal{B}$ $1$ točka, za grafično upodobitev $1$ točka.
    \end{solution}


    % 4. vprašanje
    \question
    Množica $\mathcal{C}$ je prava podmnožica množice $\mathcal{B}$, prav tako je $\mathcal{B}$ prava podmnožica množice $\mathcal{A}$.
    \begin{parts}
        \part[1]
        Grafično predstavite množico $(\mathcal{B}\setminus\mathcal{C})\cup\mathcal{A}^\complement$.
        \part[3]
        Izračunajte $\mathcal{B}\setminus\mathcal{A}$, $\mathcal{C}\cap\mathcal{A}$ in $\mathcal{C}\cup\mathcal{B}$.
    \end{parts}

    \begin{solution}[1cm]
        (a): grafična upodobitev \\
        (b): $\mathcal{B}\setminus\mathcal{A}=\emptyset$; $\mathcal{C}\cap\mathcal{A}=\mathcal{C}$; $\mathcal{C}\cup\mathcal{B}=\mathcal{B}$ \\ \\
        Pri (a) za narisan Euler-Vennov diagram $1$ točka. Pri (b) za izračun vsake izmed množic po $1$ točka.
    \end{solution}




    % \newpage

    % 5. vprašanje
    \question
    Naj bo $\mathcal{U}=\mathbb{N}_{24}$. Dane so množice $\mathcal{A}=\left\{x;\ x=2n \land 3\leq n<10\right\}$, $\mathcal{B}=\left\{x;\ x=3k-2 \land 3\leq k\leq 7\right\}$ in $\mathcal{C}=\left\{x;\ x=5m \land 1<m\leq 4\right\}$. 
    \begin{parts}
        \part[5]
        Določite naslednje množice: $\mathcal{A}\cap\mathcal{C}$, $\mathcal{C}\setminus\mathcal{A}$, $(\mathcal{A}\cup\mathcal{C})\setminus\mathcal{B}$ in $(\mathcal{A}\cup\mathcal{B})^\complement$.
        \part[1]
        Koliko elementov ima množica $\mathcal{C}\times\mathcal{B}$?
    \end{parts}

    \begin{solution}[9cm]
        $\mathcal{A}=\left\{6,8,10,12,14,16,18\right\}$, $\mathcal{B}=\left\{7,10,13,16,19\right\}$, $\mathcal{C}=\left\{10,15,20\right\}$ \\
        (a): $\mathcal{A}\cap\mathcal{C}=\left\{10\right\}$; $\mathcal{C}\setminus\mathcal{A}=\left\{15,20\right\}$; $(\mathcal{A}\cup\mathcal{C})\setminus\mathcal{B}=\left\{6,8,12,14,15,18,20\right\}$; $(\mathcal{A}\cup\mathcal{B})^\complement=\left\{1,2,3,4,5,9,11,15,17,20,21,22,23,24\right\}$ \\
        (b): $m(\mathcal{C}\times\mathcal{B})=m(\mathcal{C})\cdot m(\mathcal{B})=3\cdot 5 = 15$ \\ \\
        Pri (a) za zapis/uporabo pravilnih množic $\mathcal{A}, \mathcal{B}, \mathcal{C}$ $1$ točka; za vsako pravilno določeno množico po $1$ točka. Pri (b) za uporabo formule in izračun moči kartezičnega produkta $1$ točka.
    \end{solution}


   \newpage 
    % 6. vprašanje
    \question
    Pri uri slovenščine so se dijaki in učitelj pogovarjali o literarnih obdobjih.  Učitelj je dijake vprašal, iz katerega obdobja so njihova največkrat prebirana literarna dela.
    $20$ dijakov je odgovorilo \textit{moderna}, $7$ dijakov je reklo \textit{romantika} in $17$ \textit{realizem}. Kar $12$ jih je odgovorilo \textit{moderna} in \textit{realizem}, $4$ so rekli \textit{moderna} in \textit{romantika},
    $3$ pa \textit{romantika} in \textit{realizem}. En sam dijaki je odgovoril z vsemi tremi obdobji, dva pa nista odgovorila. 
    \begin{parts}
        \part[1]
        Narišite diagram.
        \part[2]
        Koliko dijakov je bilo pri uri slovenščine?
        \part[2]
        Koliko dijakov ima največkrat prebrano delo iz samo enega obdobja?
    \end{parts}

    \begin{solution}[8cm]
        (a): Euler-Vennov diagram z vpisanimi številkami \\
        (b): $5+1+3+3+2+11+1+2=28$ dijakov \\
        (c): $5+3+1=9$ dijakov \\ \\
        Pri (a) za pravilno narisan in izpolnjen diagram $1$ točka. Pri (b) in (c) za pravilen izračun $1$ točka, za zapis odgovora $1$ točka.
    \end{solution}




    % \newpage
        

    % 7. vprašanje
    \question[4]
    Izračunajte. $$ \left(3^2-2^3\right)^6\left(4^2-(-5)^2\right)^2+\left((-1)^9(-2)^2+(-3)^2-2^2\right)^{2024}$$
    \begin{solution}[5cm]
            $\begin{aligned}
                &(3^2-2^3)^6(4^2-(-5)^2)^2+((-1)^9(-2)^2+(-3)^2-2^2)^{2024}= \\
                    =& (9-(-8))^6(16-25)^2+(-1\cdot 4+9-4)^{2024}=\\
                    =& 1^6\cdot(-9)^2+1^{2024}= \\
                    =& 1\cdot 81+1= \\
                    =& 82
            \end{aligned} $ \\ \\
            Za upoštevanje sodih/lihih eksponentov glede na predznak $1$ točka, upoštevanje vrstnega reda operacij $1$ točka, izračun potenc $1$ in $-1$ $1$ točka, končen izračun $1$ točka.
    \end{solution}

    
    %dodatna naloga
    \bonusquestion[3]
    \textit{Dodatna naloga:} Izjavo \textit{Obstaja takšno celo število $x$, da je deljivo s številom $3$ in s številom $2$.} zapišite s simboli in jo zanikajte v simbolnem zapisu.

    \begin{solution}
        $\exists x\in\mathbb{Z}\ni:(3\mid x)\land(2\mid x)$; negacija: $\forall x\in\mathbb{Z}: (3\nmid x)\lor(2\nmid x)$ \\ \\
        Za pravilen zapis izjave s simboli $1$ točka, za pravilno negacijo kvantifikatorja $1$ točka, za pravilno negacijo logične operacije $1$ točka.
    \end{solution}


\end{questions}



\end{document}